\section{Keep linked values together when sorting}
\textbf{Recognize:} each object has several fields, but only one field
chooses the order. Sorting each array independently destroys the objects.
Prefer one record per object when you control the representation:
\begin{lstlisting}
struct Row { int key, price, quantity, id; };
vector<Row> rows;
stable_sort(rows.begin(), rows.end(),
    [](const Row& a, const Row& b) {
        return a.key < b.key;
    });
\end{lstlisting}
If the input already uses parallel arrays, sort indices once and apply the
same permutation to every array. With keys $[4,1,4]$, prices $[40,10,41]$
and quantities $[2,9,3]$, the permutation $[1,0,2]$ preserves both equal-key
records in input order. Sorting prices separately would lose that association.
\sourceimport{finals-practical}{Prelude.h}
\sourceimport{finals-practical}{LinkedSort.h}
\begin{lstlisting}
auto p = finals_practical::sorted_indices(arr1);
arr1 = finals_practical::reordered(arr1, p);
arr2 = finals_practical::reordered(arr2, p);
arr3 = finals_practical::reordered(arr3, p);
\end{lstlisting}
All arrays must have the same length. This is linked sorting, not coordinate
compression: compression replaces distinct values with ranks. For a secondary
key, compare \texttt{tie(a.key,a.price)}; for descending primary order reverse
only that comparison. Never use \texttt{<=} in a sort comparator. Do not
compare by subtraction, which can overflow. Stable input-order ties are a
choice, not the correct tie rule for every DP; strict dominance needs the
batching described in Fundraising.

\section{Recipes and inventory: validate, then commit}
\textbf{Recognize:} an operation consumes several shared resources and must
succeed completely or leave the state unchanged. Accumulate its requirements
first, check every limit, then subtract once. This pattern applies to recipes,
warehouse orders, machine resources, and bundles of tickets.

If stock is $[5,4]$, recipe A needs $(0,2),(1,1)$ and recipe B needs $(0,2)$,
order $[A,B]$ consumes $[4,1]$, leaving $[1,3]$. A second order $[A]$ fails
and must leave $[1,3]$ unchanged. Repeated dishes and repeated ingredient IDs
within a recipe both add to the requirement; a set would be incorrect.
\sourceimport{finals-practical}{AtomicInventory.h}
\begin{lstlisting}
finals_practical::Inventory inv(stock);
int accepted = 0;
for (const auto& order : orders) {
    if (!inv.consume(recipes, order)) break;
    ++accepted;
}
\end{lstlisting}
This loop answers the number of consecutive fulfillable orders, stopping at
the first failure, as in the supplied example. If the problem says to reject
an impossible order and continue, use \texttt{continue} instead of
\texttt{break}. If it lets you choose or reorder orders to maximize acceptance,
this simulation does not solve the optimization problem.

\textbf{Bounds:} use 64-bit stock and quantities. Each is nonnegative;
check \texttt{q > stock[id]-need[id]} before adding, so even an oversized
request cannot overflow. All recipe and ingredient IDs must be valid; convert
one-based input exactly once. Zero quantities and empty orders are supported.
Negative quantities would invalidate the early-failure argument.

The touched-index list clears only ingredients examined in this order, rather
than allocating and scanning an $N$-entry array every time. Initialization is
$O(N)$; each call costs its dish count plus expanded ingredient entries and
touched indices. Repeated large recipes still cost work per occurrence.
Orders may be read and processed one at a time to avoid storing them all.
When an early order fails, still read remaining input if a surrounding
multi-test-case format requires it.

\section{Offline queries: activate values in threshold order}
\textbf{Recognize:} a fixed array, many interval queries, and each query
counts values at least a threshold. Queries can be answered out of input
order, then restored using their original indices. Sort values and queries
by decreasing threshold. Before answering a query, activate every array
position whose value is at least its threshold. A Fenwick tree stores one
at active positions, so a range sum gives the count.
\sourceimport{finals-practical}{OfflineThreshold.h}
For $a=[5,2,7,4]$ and query $[1,4)$ with threshold 5, positions 0 and 2
are active; only position 2 lies in the interval, so the answer is 1.
At threshold 4, activate position 3 too and the answer becomes 2.
Equal values must be activated before an inclusive query. For strictly
greater values, change the activation comparison to \texttt{>}.

This preserves the original array indices; sorting the values and then
querying their sorted positions would answer a different question. All
ranges must satisfy $0\le l\le r\le n$. The array is static: interleaved
updates require a different method. After Euler flattening, this same
helper answers threshold queries on subtrees. The output order matches
the query input order even though evaluation order changes.
